{\footnotesize
\begin{tabularx}{\textwidth}{>{\fontSemiBold\color{BrandNavy}\hsize=0.535\hsize\sloppy\arraybackslash}X >{\hsize=1.206\hsize\sloppy\arraybackslash}X >{\hsize=1.259\hsize\sloppy\arraybackslash}X}
\arrayrulecolor{GithubTableBorder}\toprule
\rowcolor{BrandTableHeader}
\textbf{Seção / Dimensão} & \textbf{Versão 1.7 (Antes da Melhoria)} & \textbf{Versão 2.0 (Após a Melhoria)} \\
\midrule
\arrayrulecolor{BrandGrayRule}
Seção 4.1: Escopo e Cronograma & Cronograma empírico sem divisão formal de entregas e sem responsáveis nominais. & EAP em 3 níveis, 3 Macroentregas definidas e responsáveis nominais por Sprint. \\
\midrule
Seção 4.1: Tratamento de Desvios & Sem procedimento documentado para acompanhamento e contingência de atrasos. & Protocolo claro de 48h no GitHub Projects com pareamento técnico ou fatiamento MoSCoW. \\
\midrule
Seção 4.2: Matriz de Riscos & Registro de riscos resumido, sem severidade quantitativa (P×I) e sem ações reativas. & Risk Register com 7 riscos reais quantificados, mitigação preventiva e contingência. \\
\midrule
Seção 4.5: Métricas de Projeto & Ausência de tabela formal de indicadores, metas e ações corretivas de desvio. & 6 métricas objetivas (Lead Time, WIP, testes, taxa de defeitos, escopo e custos). \\
\arrayrulecolor{GithubTableBorder}\bottomrule
\end{tabularx}%
}
